\documentclass[11pt]{article}
\usepackage{preamble}
\setlength{\parskip}{4pt}

\begin{document}

\daybanner{AP Statistics \textperiodcentered\ Topic 4.3 (CED 2.4) \textperiodcentered\ B: 2026-10-22 \textperiodcentered\ E: 2026-11-06 \textperiodcentered\ TEACHER KEY}{Probability or Estimate?}

\sectionbanner{CED TETHER}

{\small
\begin{itemize}[leftmargin=1.5em,itemsep=2pt,topsep=2pt]
  \item \textbf{Skill 3.A} --- Determine relative frequencies, proportions, or probabilities using simulation or calculations.
  \item \textbf{Skill 4.B} --- Interpret statistical calculations and findings to assign meaning or assess a claim.
  \item \textbf{EU VAR-4} --- The likelihood of a random event can be quantified.
  \item \textbf{LO VAR-4.A} --- Calculate probabilities for events and their complements. [Skill 3.A]
  \item \textbf{EK VAR-4.A.1} --- The sample space of a random process is the set of all possible non-overlapping outcomes.
  \item \textbf{EK VAR-4.A.2} --- If all outcomes in the sample space are equally likely, then the probability an event E will occur is defined as the fraction: (number of outcomes in event E) / (total number of outcomes in sample space).
  \item \textbf{EK VAR-4.A.3} --- The probability of an event is a number between 0 and 1, inclusive.
  \item \textbf{EK VAR-4.A.4} --- The probability of the complement of an event E, E' or E\textasciicircum{}C (i.e., not E), is equal to 1 - P(E).
  \item \textbf{LO VAR-4.B} --- Interpret probabilities for events. [Skill 4.B]
  \item \textbf{EK VAR-4.B.1} --- Probabilities of events in repeatable situations can be interpreted as the relative frequency with which the event will occur in the long run.
\end{itemize}
}
\textbf{Essential question:} How should a reader connect a theoretical probability to what happened in a finite set of trials?\par
\textbf{DOK-3 skill:} 4.B \textperiodcentered\ \textbf{FRQ pattern:} {\small\texttt{theoretical-\allowbreak{}vs-\allowbreak{}empirical-\allowbreak{}which-\allowbreak{}is-\allowbreak{}the-\allowbreak{}probability}}\par
\textbf{DOK rationale:} Part (c) coordinates an equally likely sample-space model with an observed relative frequency, adjudicates the roles of two competing numerical claims, and bounds when the empirical estimate deserves trust.\par

\frameworkphaseheader{Do Now (first take) $\rightarrow$ Explore (video + follow-along) $\rightarrow$ Exit (finish + turn in)}{1 $\rightarrow$ 3}{45}{%
\begin{itemize}[leftmargin=1.5em,itemsep=2pt,topsep=2pt]
  \item Board slide up at the bell. Ask students to commit to one number without confirming either claim.
  \item Begin the combined 4.3--4.5 video follow-along at minute 5.
  \item Return to the board slide at minute 33. Circulate and require students to label each number by its source.
  \item Collect at the door. Score part (c) only, E/P/I.
\end{itemize}
}{%
\begin{itemize}[leftmargin=1.5em,itemsep=2pt,topsep=2pt]
  \item Choose 0.375 or 0.45 and cite either the spinner design or the observed plays.
  \item Complete the follow-along about sample spaces, long-run interpretations, intersections, and conditional probability.
  \item Finish (a)--(c), revise the first take in the exit line, and turn in the sheet.
\end{itemize}
}{%
\begin{itemize}[leftmargin=1.5em,itemsep=2pt,topsep=2pt]
  \item Which number comes from possible outcomes, and which comes from recorded outcomes?
  \item Does a difference of 7.5 percentage points prove that either calculation is wrong?
  \item What exactly does the law of large numbers predict as trials accumulate?
  \item If the spinner favored one sector, which number would stop representing the actual game?
\end{itemize}
}{Solo teacher. Circulate ~5 pairs. Ask for the source and role of each number; do not accept `the data are wrong' or `the model always wins' without a condition.}

\begin{callout}{\IconWarn\ WATCH FOR}{calloutred}
\begin{itemize}[leftmargin=1.5em,itemsep=2pt,topsep=2pt]
  \item Treating the 40-play relative frequency as an exact, permanent probability.
  \item Treating the law of large numbers as a fixed-trial guarantee.
  \item Using the equal-sector count even if the physical spinner or spinning method is not fair.
\end{itemize}
\end{callout}

\sectionbanner{SCORING PART (c) --- E / P / I}

\textbf{Expected elements}
\begin{itemize}[leftmargin=1.5em,itemsep=2pt,topsep=2pt]
  \item \textbf{distinguishes-theoretical-empirical} (required) --- Identifies 3/8 as the theoretical model probability and 18/40 as the empirical estimate
  \item \textbf{uses-both-sources} (required) --- Cites three winning outcomes among eight equal sectors and 18 wins among 40 observed plays
  \item \textbf{invokes-long-run} (required) --- Uses the law of large numbers to explain chance variation now and convergence over many comparable trials
  \item \textbf{bounds-trust} (required) --- Checks both the number of trials and whether the spinner and procedure produce equally likely outcomes
\end{itemize}

{\small\renewcommand{\arraystretch}{1.25}
\noindent\begin{tabular}{|p{0.5in}|p{5.9in}|}
\hline
\textbf{E} & Correctly assigns the two numbers their roles, cites both sources, explains the difference with the law of large numbers, and gives both a trial-count and a fairness limitation. \\ \hline
\textbf{P} & Correctly distinguishes probability from estimate and cites the data, but gives an incomplete long-run explanation or only one limitation on trust. \\ \hline
\textbf{I} & Treats 0.45 as the guaranteed probability solely because it was observed, dismisses the data because it differs from 3/8, or says the law of large numbers forces exactly 37.5 percent wins in a fixed number of plays. \\ \hline
\end{tabular}}

\textbf{Common mistakes}
\begin{itemize}[leftmargin=1.5em,itemsep=2pt,topsep=2pt]
  \item Calling $18/40$ the exact probability even under the stated equal-sector model
  \item Saying the two values differ because one calculation must be wrong
  \item Claiming the law of large numbers guarantees an exact match after a specific number of plays
  \item Ignoring that unequal sectors or a biased spinning method would invalidate the equally likely model
\end{itemize}

\clearpage
\focusbanner{TODAY'S PROBLEM --- annotated}

Suppose a game uses a spinner with eight equal sectors numbered 1 through 8. A player wins when the spinner lands on 1, 2, or 3. The dotplot shows the results of 40 observed plays.\par\smallskip \textbf{Dani:} ``The probability of winning is $3/8=0.375$ because three of the eight equally likely results win.''\par \textbf{Micah:} ``The probability is $18/40=0.45$ because 18 of the 40 observed plays were wins.''\par
\begin{center}
\begin{tikzpicture}
\begin{axis}[
  width=0.61\linewidth,
  height=1.15in,
  ymin=0, ymax=9, ytick=\empty, axis y line=none, axis x line=bottom,
  xlabel={Spinner result (40 plays)}, tick label style={font=\small}, label style={font=\small}
]
\addplot[only marks, mark=*, mark size=1.7pt, color=dayblue] coordinates {
  (1,1)
  (1,2)
  (1,3)
  (1,4)
  (1,5)
  (1,6)
  (1,7)
  (1,8)
  (2,1)
  (2,2)
  (2,3)
  (2,4)
  (2,5)
  (2,6)
  (3,1)
  (3,2)
  (3,3)
  (3,4)
  (4,1)
  (4,2)
  (4,3)
  (4,4)
  (4,5)
  (5,1)
  (5,2)
  (5,3)
  (6,1)
  (6,2)
  (6,3)
  (6,4)
  (6,5)
  (6,6)
  (7,1)
  (7,2)
  (7,3)
  (7,4)
  (8,1)
  (8,2)
  (8,3)
  (8,4)
};
\end{axis}
\end{tikzpicture}
\end{center}
\begin{firsttakebox}{FIRST TAKE (5 min)}
Which number would you use to describe the chance of winning the next play: $0.375$ or $0.45$? Give one reason.\par
\answer{Not graded. Expect both choices; the revision should distinguish a model probability from a data-based estimate.}
\end{firsttakebox}

\textit{Rules callout ``MODEL, DATA, AND THE LONG RUN (keep on the page)'' is printed on the student sheet.}\par\medskip

\dokbadge{1}~\textbf{(a)} List the sample space and identify the outcomes in the event `win.'\par
\answer{The sample space is $S=\{1,2,3,4,5,6,7,8\}$. The win event is $W=\{1,2,3\}$.}

\dokbadge{2}~\textbf{(b)} Use the dotplot to verify the observed relative frequency of a win. Compute the theoretical probability under the equal-sector model, then give the difference in percentage points.\par
\answer{The dotplot has 8 ones, 6 twos, and 4 threes, so there are $18$ wins and the observed relative frequency is $18/40=9/20=0.45$. With eight equally likely sectors and three winning results, $P(W)=3/8=0.375$. The difference is $0.45-0.375=0.075$, or $7.5$ percentage points.}

\dokbadge{3}~\textbf{(c)} Under the stated equal-sector model, decide which number is the theoretical probability and which is an empirical estimate. Cite both the sample-space feature and the 40-play result, use the law of large numbers to explain why the numbers can differ, and state two features of the trial process you would check before trusting the observed proportion.\par
\answer{Dani's $3/8=0.375$ is the theoretical probability under the model because 3 of 8 equally likely sectors win. Micah's $18/40=0.45$ is an empirical estimate from this particular set of 40 plays. Random variation can make a short-run proportion differ from the theoretical probability; the law of large numbers says the observed proportion tends to approach the true probability as the number of independent, identically conducted plays grows, not that 40 plays must match exactly. Before trusting $0.45$, check whether there were enough plays and whether the spinner and spinning method made the eight outcomes equally likely. If the game was not fair, $3/8$ would not be its actual probability; if only 40 fair plays were observed, $0.45$ may simply be a variable estimate.}

\begin{summaryexitbox}{EXIT}
One thing the video changed about my first take: \hrulefill
\end{summaryexitbox}

\end{document}
